
\documentclass[conference]{IEEEtran}
\IEEEoverridecommandlockouts

\usepackage{cite}
\usepackage{amsmath,amssymb,amsfonts}
\usepackage{algorithmic}
\usepackage{graphicx}
\usepackage{textcomp}
\usepackage{xcolor}
\usepackage{hyperref}
\usepackage{booktabs}
\usepackage{caption}
\usepackage{xcolor}
\usepackage{verbatim}
\hypersetup{
  colorlinks   = true,
  linkcolor    = blue,
  citecolor    = blue,
  urlcolor     = blue,
}

% --- Useful macros ---
\newcommand{\E}{\mathbb{E}}
\newcommand{\R}{\mathbb{R}}
\newcommand{\D}{\mathcal{D}}
\newcommand{\Dir}{\mathrm{Dir}}
\newcommand{\OrbitALIF}{\textsc{OrbitALIF}}
\newcommand{\FQS}{5QS}
\newcommand{\DFRB}{\textsc{DFRB}}
\newcommand{\SHAM}{\textsc{SHAM}}
\newcommand{\AGFM}{\textsc{AGFM}}
\newcommand{\SSHB}{\textsc{SSHB}}

\graphicspath{{figures/}}

% --- Tighten float spacing to fit within the 6-page IEEE limit ---
\setlength{\textfloatsep}{6pt plus 1pt minus 1pt}
\setlength{\floatsep}{6pt plus 1pt minus 1pt}
\setlength{\intextsep}{6pt plus 1pt minus 1pt}
\setlength{\dblfloatsep}{6pt plus 1pt minus 1pt}
\setlength{\dbltextfloatsep}{6pt plus 1pt minus 1pt}
\setlength{\abovecaptionskip}{3pt}
\setlength{\belowcaptionskip}{0pt}

\begin{document}

% --- Title block ---
\title{ORBITALIF: An Efficient Spiking Federated Learning Framework for Onboard Cloud Removal\\
\thanks{This work has been supported by the National Nature Science Foundation of China under Grant 62571331. (\textit{Corresponding author: Yijie Mao})}
}

\author{
\IEEEauthorblockN{Bohan Zhang, Chenyu Xu, Yijie Mao}
\IEEEauthorblockA{
School of Information Science and Technology, ShanghaiTech University, Shanghai, China\\
Email: \{xuchy2024, maoyj\}@shanghaitech.edu.cn}
}

\maketitle

% =====================================================================
% Abstract
% =====================================================================
\begin{abstract}
Optical low-earth-orbit (LEO) satellites enable high-resolution, large-scale Earth observation for applications such as disaster monitoring and environmental surveillance. However, cloud coverage often obscures the Earth’s surface, and conventional cloud-removal pipelines that download cloudy images to ground stations for processing suffer from limited contact windows, constrained satellite-to-ground bandwidth, and high latency. 
%%%
In this work, we propose \OrbitALIF{}, a novel energy-efficient framework that enables both onboard training and onboard inference for optical cloud removal. \OrbitALIF{} performs local training on each satellite, exchanges only model parameters through inter-satellite links, and employs a lightweight spiking neural network (SNN) for energy-efficient onboard inference.
%%%
To meet satellite power constraints, we design a compact $2.30$M-parameter SNN backbone with sparse event-driven computation and multi-domain feature modeling for accurate cloud removal. To address heterogeneous observations across satellites, \OrbitALIF{} adopts a federated batch normalization (FedBN)-based aggregation strategy that shares common model weights while keeping satellite-specific normalization statistics local to each satellite.
%%%
On the public CUHK-CR1 optical remote sensing cloud-removal dataset, \OrbitALIF{} achieves a peak signal-to-noise ratio (PSNR) of $25.37$ dB while reducing per-image inference energy by $98.6\%$ compared with an artificial neural network (ANN) baseline. Despite its lightweight design, it achieves a PSNR only $0.81$ dB lower than the existing best cloud-removal model and demonstrates strong robustness towards heterogeneous data distributions.

%Optical low-earth-orbit (LEO) satellites collect large volumes of Earth imagery, but clouds often hide the surface and make the captured data unusable for time-sensitive applications such as disaster response. A conventional solution is to downlink cloudy images to ground stations and remove clouds on the ground. This pipeline, however, suffers from limited contact windows, constrained satellite-to-ground bandwidth, and hours-level latency. A more responsive solution is to perform cloud removal directly onboard the satellite constellation.

%Onboard cloud removal is difficult for two reasons. First, each satellite has only a limited power budget, making large artificial neural networks expensive to deploy. Second, satellites observe different regions and cloud conditions, but cannot upload all raw images to a central server for training. This requires a learning framework that is both energy-efficient and communication-aware.

%We present \OrbitALIF{}, the first end-to-end framework that performs both onboard training and onboard inference for optical cloud removal.
%A $2.30$\,M-parameter spiking neural network backbone is coupled with
%a FedBN-on-SNN aggregation protocol that runs over intra-plane ring
%AllReduce and inter-plane Gossip mixing. 

%On CUHK-CR1, \OrbitALIF{} reaches $25.37$\,dB PSNR while reducing per-image inference energy from $20.75$\,mJ for an ANN baseline to $0.287$\,mJ on neuromorphic hardware, a $72.3$-fold reduction. It also stays within $0.81$\,dB of the strongest published cloud-removal model on this benchmark at a fraction of their parameter count. Under Dirichlet source heterogeneity of different $\alpha$, the architectural proves its robust.
\end{abstract}

\begin{IEEEkeywords}
Spiking neural networks, optical cloud removal, decentralized
federated learning, onboard satellite computing.
\end{IEEEkeywords}

% =====================================================================
% Section I: Introduction
% =====================================================================
\section{Introduction}\label{sec:intro}

Low-earth-orbit (LEO) satellites have become an essential
infrastructure for Earth observation, supporting time-sensitive
applications such as disaster response~\cite{ref:earthobs2},
environmental monitoring~\cite{ref:earthobs1}, and large-scale
remote-sensing analytics.
Cloud removal, which aims to recover the origin surface from a
cloud-contaminated observation, is therefore a key preprocessing step
for remote-sensing tasks.

Existing cloud-removal methods have mainly been developed for
ground-side processing. Early approaches rely on multi-temporal
compositing or synthetic-aperture-radar (SAR) multi-sensor fusion.
Meraner~\emph{et~al.}~\cite{ref:meraner2020} showed that a deep
residual network conditioned on co-registered SAR data effectively
reconstructs cloud-occluded regions, setting a strong baseline for
subsequent fusion methods. Ebel~\emph{et~al.}~\cite{ref:ebel2021multisensor}
extended this line of work to global, all-season settings, but their
pipeline requires accurate registration and is therefore less suitable
for time-critical scenarios where only a single cloudy snapshot is
available. To overcome this limitation, single-image deep-learning
methods have been proposed: GAN- and CNN-based models such as
SpA-GAN, AMGAN-CR, and MemoryNet directly learn a mapping from a
cloudy tile to its cloud-free counterpart, while more recent
transformer- and diffusion-based models further improve restoration
fidelity through long-range attention and iterative
denoising~\cite{ref:cuhkcr}. Although these models provide strong
restoration quality, their computational cost remains prohibitive for
onboard satellite deployment.


A natural way to reduce response latency is to move cloud removal
onto the satellite constellation, which avoids downlinking large
volumes of cloud-contaminated imagery. However, onboard cloud removal
must satisfy several coupled physical constraints. First, small and
medium LEO satellites have tight power budgets, typically on the
order of $100$--$300$\,W for the whole bus~\cite{ref:spacecpn1}.
Second, satellite-to-ground and inter-satellite links have limited
bandwidth and intermittent contact windows~\cite{ref:spacecpn1}.
Third, constellations have asymmetric connectivity: intra-plane links
are relatively stable, whereas inter-plane links are available only
under restricted orbital visibility conditions~\cite{ref:walker1984}.
Fourth, different satellites observe different geographic regions,
cloud climatologies, illumination conditions, and sensor statistics,
causing source-level feature heterogeneity.

Federated learning (FL) provides a solution by letting each satellite
train locally and share only model parameters~\cite{ref:fedavg}.
Early satellite-FL work applied FedAvg directly to satellite--ground
scenarios with visibility-aware scheduling. Subsequent work moved
aggregation into the constellation: FedMega~\cite{ref:fedmega2024}
introduced ring all-reduce over intra-orbit ISLs to cut ground-link
dependence, and DFedSat~\cite{ref:dfedsat} further removed the ground
server by combining intra-orbit reduce with inter-orbit gossip over
dynamic ISLs. However, prior efforts have focused on classification
and detection tasks; dense pixel-level cloud removal, which demands
per-pixel regression under physical non-IID, remains unaddressed.

% Dense cloud removal is more difficult because the model must reconstruct every pixel while preserving local texture, global color consistency,and cloud-boundary structure.

Energy-efficient neural restoration is another necessary component
for onboard deployment. Spiking neural networks (SNNs) are attractive
because their event-driven computation can skip synaptic operations
when no spike is emitted~\cite{ref:loihi}. Recent low-level-vision
work has explored SNNs for image restoration tasks such as
deraining~\cite{ref:esdnet}. However, conventional binary spikes
provide limited activation expressiveness, which is problematic for
dense pixel-level regression. Aggressively reducing computation
degrades restoration quality, while preserving quality often requires
deeper networks that weaken the energy advantage. Therefore, onboard
cloud removal requires a co-design of the neuron activation,
restoration backbone, and constellation-level aggregation protocol.

Despite progress in cloud removal, satellite FL, and SNN-based
efficient inference, onboard federated cloud removal remains largely
unaddressed. The task involves three coupled challenges. First, it is
a dense regression problem rather than a
classification task. Second, cloud artifacts are multi-scale and
multi-domain, ranging from thin high-frequency boundaries to thick
low-frequency occlusions. Third, the non-IID distribution arises from
physical source shifts across orbital planes rather than only from
label imbalance.  To the best of our knowledge,
there is no prior end-to-end framework that jointly supports onboard
training and onboard inference for  cloud removal over a LEO
satellite constellation.

%
The main contributions of this paper are summarized as follows:

\begin{itemize}

\item We propose \OrbitALIF{}, an onboard SNN-based federated
cloud-removal framework for LEO satellite constellations. To the best
of our knowledge, this is the first work to go beyond classification
tasks and combine onboard SNN inference with cross-orbit federated
training for a low-level vision task such as cloud removal. 

\item We design a lightweight SNN backbone that analyses
cloud-contaminated imagery jointly across temporal, spectral, and
spatial domains to extract multi-dimensional feature representations.
The backbone comprises four dedicated modules and contains only
$2.30$\,M parameters.

\item We evaluate \OrbitALIF{} on the public CUHK-CR cloud-removal
benchmark over a $5{\times}10$ Walker-Star constellation. Under
intra-plane ring all-reduce and inter-plane chain-Gossip with
plane-local BN, the federated \OrbitALIF{} attains restoration
quality comparable to an ANN U-Net of the same architecture while
reducing per-image inference energy by $72.3\times$. Under
source-level Dirichlet heterogeneity at $\alpha\!\in\!\{0.1,0.01\}$,
the framework remains robust to extreme non-IID conditions.
\end{itemize}



% Section II: System Model
% =====================================================================
\section{System Model}\label{sec:sysmodel}

\subsection{Satellite Communication Model}\label{sec:sysmodel-setup}

In this work, we consider a constellation consisting of $N$ orbital planes, each containing $K$ satellites. We denote the $j$-th satellite within the $i$-th orbital plane by $s_{i,j}$. Satellites are assumed to communicate with one another via inter-satellite links (ISLs). Specifically, two decision metrics are introduced to determine whether a pair of satellites can establish communication: line-of-sight (LoS) visibility and Doppler reliability, which are elaborated as follows:
\begin{itemize}
    \item \emph{LoS visibility.} Consider an arbitrary satellite pair
    $s_{i_1,j_1},s_{i_2,j_2}$ (hereafter $s_1,s_2$) at altitudes
    $h_1,h_2$. Each satellite's horizon half-angle is
    $\rho_k = \arccos\!\bigl(r_E / (r_E + h_k)\bigr)$, $k\in\{1,2\}$,
    where $r_E$ is the Earth radius. Letting
    $(\theta_k^{v},\theta_k^{h})$ denote the latitude and longitude
    of $s_k$, the geocentric angle between the two satellites is
    \begin{equation*}
        \varphi = \arccos\!\bigl[\sin\theta_{1}^{v}\sin\theta_{2}^{v}
        + \cos\theta_{1}^{v}\cos\theta_{2}^{v}\cos(\theta_{1}^{h} - \theta_{2}^{h})\bigr].
    \end{equation*}
    The two satellites are mutually LoS-visible if and only if
    $\varphi \leq \rho_1 + \rho_2$.
\item \emph{Doppler reliability.} Even with a clear line of sight,
    a high relative velocity can render an ISL unusable. The Doppler
    shift between $s_1,s_2$ is $f_d = \psi f_c / c$, where $\psi$ is
    the relative radial speed, $f_c$ the carrier frequency, and $c$
    the speed of light. The two satellites are Doppler-reliable if
    $f_d \leq f_{\max}$ for a tolerable threshold $f_{\max}$.
\end{itemize}

Once both metrics are satisfied, the satellite pair is deemed a credible pair. Upon the formation of a credible pair, information can be successfully transmitted between the two satellites.

\subsection{Satellite Federated Learning Model}

Assume each satellite $s_{i,j}$ holds a local dataset $\mathcal{D}_{i,j}$, all satellites collaboratively learn a global model by solving the following empirical risk minimization (ERM) problem:
\begin{equation*}
\begin{aligned}
\min_{\mathbf{x} \in \mathbb{R}^d} \quad & f(\mathbf{x}) = \frac{1}{NK} \sum_{i=1}^{N} \sum_{j=1}^{K} f_{i,j}(\mathbf{x}), \\
& f_{i,j}(\mathbf{x}) = \frac{1}{|\mathcal{D}_{i,j}|} \sum_{\mathbf{z} \in \mathcal{D}_{i,j}} \ell(\mathbf{x}; \mathbf{z}),
\end{aligned}
\end{equation*}
where $\mathbf{x}$ denotes the model parameters and
$\ell(\cdot;\mathbf{z})$ is the loss on sample $\mathbf{z}$. We
partition the parameter vector into a \emph{shared} component
$\mathbf{x}^{\mathrm{sh}}$ (convolutional weights, learnable scalars,
frequency-MLP weights) and a \emph{local} component
$\mathbf{x}^{\mathrm{bn}}$ (batch-normalization affine parameters and
running statistics, including the threshold-dependent BN of the
spiking blocks). Each global round $t$ proceeds in four steps.

\emph{Step 1 -- Model distribution.} At the start of round $t$, the
$K$ satellites in plane $i$ share an identical plane model
$\mathbf{x}_i^t$, established at the end of round $t-1$ and
disseminated within the plane via the stable intra-plane ring.

\emph{Step 2 -- Local training.} Each satellite $s_{i,j}$ performs
local stochastic gradient descent (SGD) on $\mathcal{D}_{i,j}$ for
$E$ local steps,
\begin{equation*}
\mathbf{x}_{i,j}^{t+\frac{1}{2}} = \mathbf{x}_i^t - \eta \sum_{e=0}^{E-1} \nabla F_{i,j}\bigl(\mathbf{x}_{i,j}^{t,e}\bigr),
\end{equation*}
where $\eta$ is the learning rate, $\nabla F_{i,j}$ a stochastic
gradient, and $\mathbf{x}_{i,j}^{t,e}$ the model after $e$ local
steps.

\emph{Step 3 -- Intra-plane aggregation (ring all-reduce).} Because
satellites within the same orbital plane move in the same direction
at nearly constant velocities, the intra-plane ISLs form a stable
logical ring. After $\mathcal{O}(K)$ communication rounds, every
satellite obtains the exact average of all local models within the
plane,
\begin{equation*}
\mathbf{x}_i^{t+\frac{1}{2}} = \frac{1}{K}\sum_{j=1}^{K} \mathbf{x}_{i,j}^{t+\frac{1}{2}}, \quad \forall i\in\{1,\dots,N\}.
\end{equation*}
Same-plane satellites observe a similar climatology slice, so this
average is applied to \emph{all} parameters, including
$\mathbf{x}^{\mathrm{bn}}$.

\emph{Step 4 -- Inter-plane aggregation (chain-Gossip with
plane-local BN).} Inter-plane ISLs are unstable owing to varying
inclinations, altitudes, and relative velocities, making
single-round global aggregation impractical. Instead, each plane
$i$ communicates only with its immediate neighbours on a pre-defined
connectivity graph $\mathcal{G} = (\mathcal{V},\mathcal{E})$, where
$\mathcal{V} = \{1,\dots,N\}$ and an edge $e_{i,i'}\in\mathcal{E}$
exists only if at least one credible pair (satisfying both LoS and
Doppler constraints of Section~\ref{sec:sysmodel-setup}) links
orbits $i$ and $i'$. Letting
$\mathcal{N}(i) = \{i' \mid e_{i,i'}\in\mathcal{E}\}$ denote the
neighbour set, the shared parameters undergo an equal-weight average
over plane $i$ together with its neighbours,
\begin{equation*}
\mathbf{x}_i^{\mathrm{sh},\,t+1} = \frac{1}{|\mathcal{N}(i)|+1}\Biggl(\mathbf{x}_i^{\mathrm{sh},\,t+\frac{1}{2}} + \sum_{i'\in\mathcal{N}(i)}\mathbf{x}_{i'}^{\mathrm{sh},\,t+\frac{1}{2}}\Biggr),
\end{equation*}
while the BN component is excluded from the inter-plane mean and
retained from the plane's own pre-aggregation snapshot,
$\mathbf{x}_i^{\mathrm{bn},\,t+1} = \mathbf{x}_i^{\mathrm{bn},\,t+\frac{1}{2}}$,
following the FedBN principle. Through repeated rounds, shared
parameters diffuse across the constellation and reach consensus,
while normalization statistics remain adapted to each plane's
climatology. This avoids the high communication overhead of global
aggregation and is robust to the intermittent nature of inter-plane
ISLs.


% =====================================================================
% Section III: The ORBITALIF Framework
% =====================================================================
\section{The \OrbitALIF{} Framework}\label{sec:framework}

\begin{figure*}[!t]
  \centering
  \includegraphics[width=0.92\textwidth]{fig_archi.png}
  \caption{\OrbitALIF{} overview. \textbf{Left:} single-satellite
  U-shaped backbone ($T\!=\!4$, $2.30$\,M params). \textbf{Top-right:}
  one round on satellite $s_{1,3}$ -- local SGD (S1), intra-plane
  mean over $K\!=\!10$ sats incl.\ BN (S2), inter-plane chain-Gossip
  on shared (non-BN) parameters (S3), broadcast back (S4).
  \textbf{Bottom:} chain-Gossip on $N\!=\!5$ planes; BN stays
  plane-local (FedBN).}
  \label{fig:archi}
\end{figure*}

% ---------------------------------------------------------------------
\subsection{Design Principles and Architectural Overview}
\label{sec:framework-overview}
% ---------------------------------------------------------------------

We adopt an SNN model as the foundation; its principal architectural
innovation is multi-domain decomposition: cloud-contaminated images
behave very differently across the pixel, Fourier, and time domains.
Figure~\ref{fig:archi} summarises the resulting backbone, the
per-round protocol on a single satellite, and the inter-plane gossip
mixing rule discussed below.

%); \textbf{P2}~identity-preserving modulation (every
%block is initialized as identity so training is stable from the start);
%\textbf{P3}~progressive coupling (cross-branch gates start closed and
%open only when training finds them useful).


\noindent\textbf{Multi-domain decomposition.} Cloud artefacts behave
differently across representations: smooth blobs in the pixel
domain, mid-band concentration in the Fourier amplitude spectrum,
and slow variation along the spike-time axis. Each module therefore
decomposes its input in its native domain and fuses the components
downstream.


%\textbf{P2~(Identity-preserving modulation).} Every modulation block
%takes the form $\mathbf{x}\mapsto\mathbf{x}+\sigma(\lambda_s)
%(\Phi(\mathbf{x})-\mathbf{x})$ with $\lambda_s$ initialized to zero,
%so each block acts as an identity at training start. Uncorrupted image
%regions pass through unmodified, and a new module cannot destabilize a
%previously-trained checkpoint.

%\textbf{P3~(Progressive coupling).} Cross-branch interactions are gated
%by scalars initialized to zero. Single-branch features mature first;
%inter-branch coupling opens gradually during training.

The backbone is a three-stage U-Net. The input cloudy image
$\mathbf{c}\in\mathbb{R}^{3\times H\times W}$, where $H$ and $W$ are
the tile height and width in pixels, is patch-embedded and replicated
along a synthetic time axis of $T\!=\!4$ steps. Stage~1 (highest
resolution) uses the spectro-temporal spike hyper-block (\SSHB{})
that processes spectral and spatial information jointly in both
encoder and decoder; stages~2 and~3 stack dual-frequency residual
blocks (\DFRB{}) with embedded spectral--spatial hybrid attention
(\SHAM{}). Each skip connection passes through an adaptive gated
fusion module (\AGFM{}) before entering the decoder.
The final prediction is
\begin{equation}
    \hat{\mathbf{y}} \;=\; \mathbf{c} \;+\;
    \mathrm{Conv}_{3\times3}\!\Bigl(
    \tfrac{1}{T}\textstyle\sum_{t=1}^{T}\mathbf{s}[t]\Bigr),
    \label{eq:globalresidual}
\end{equation}
where $\hat{\mathbf{y}}\in\mathbb{R}^{3\times H\times W}$ is the
restored cloud-free image, $\mathbf{s}[t]\in\mathbb{R}^{C\times H
\times W}$ is the spike feature map at time step $t$ with $C$ channels,
and $\mathrm{Conv}_{3\times3}$ maps the temporally averaged features to
3 RGB channels. The global residual lets the network predict only the
cloud correction rather than the full image. 

%The onboard variant uses base channel width $C\!=\!24$ and has approximately $2.30$\,M parameters. Table~\ref{tab:contribution} maps each module to the principles it instantiates.

%%%
%\begin{table}[!t]
%\centering
%\caption{Module--principle matrix. MDD: multi-domain decomposition;
%IPM: identity-preserving modulation; PC: progressive coupling;
%GRF: global receptive field; SNN: spike-compatibility.}
%\label{tab:contribution}
%\setlength{\tabcolsep}{4pt}
%\renewcommand{\arraystretch}{0.95}
%\footnotesize
%\begin{tabular}{lccccc}
%\toprule
%Component & MDD & IPM & PC & GRF & SNN \\
%\midrule
%\DFRB{}               & spat./val.   & ---          & \checkmark   & ---          & \checkmark \\
%\SHAM{}               & freq./time   & \checkmark   & \checkmark   & \checkmark   & \checkmark \\
%\AGFM{}               & enc./dec.    & ---          & \checkmark   & ---          & \checkmark \\
%\SSHB{}               & spat./time   & \checkmark   & ---          & \checkmark   & \checkmark \\
%\midrule
%\textbf{\OrbitALIF{}} & \textbf{all} & \textbf{\checkmark} & \textbf{\checkmark} & \textbf{\checkmark} & \textbf{\checkmark} \\
%\bottomrule
%\end{tabular}
%\end{table}

%%%
% ---------------------------------------------------------------------
\subsection{Neuron Scale: The \FQS-LIF{} Neuron}
\label{sec:framework-neuron}
% ---------------------------------------------------------------------

A binary LIF spike carries only one bit of information per step;
matching ANN accuracy therefore requires either many time steps,
which is too costly for onboard deployment, or deeper SNN restoration
networks that erode the energy advantage~\cite{ref:esdnet}. We
instead extend the output alphabet to five levels
$\mathcal{S} = \{0,\tfrac{1}{4},\tfrac{1}{2},\tfrac{3}{4},1\}$,
packing more information into each spike so that comparable quality
is achieved at $T\!=\!4$. We call this neuron \FQS-LIF{}. Its
membrane update and spike output are
\begin{equation}
    u[t] = \bigl(u[t{-}1] - s[t{-}1]\bigr)\eta + \mathbf{W}\mathbf{x}[t], \quad
    s[t] = \tfrac{1}{4}\,\mathrm{round}\!\bigl(\mathrm{clamp}(u[t],0,4)\bigr),
    \label{eq:5qs}
\end{equation}
where $u[t]\in\mathbb{R}$ is the membrane potential at step $t$,
$s[t{-}1]\in\mathcal{S}$ the spike output from the previous step,
$\eta\in(0,1)$ the leak factor, $\mathbf{W}$ the synaptic weight
matrix, and $\mathbf{x}[t]$ the pre-synaptic input. The soft reset
subtracts $s[t{-}1]$ rather than a fixed threshold, carrying residual
charge into the next step. Synaptic operations are skipped when
$s[t]\!=\!0$, preserving event-driven sparsity.

% ---------------------------------------------------------------------
\subsection{Block Scale: DFRB and SHAM}
\label{sec:framework-block}
% ---------------------------------------------------------------------

\noindent\textbf{Dual-Frequency Residual Block (\DFRB{}).}
Cloud artefacts span different frequencies: thin edges are
high-frequency, thick haze is low-frequency. Uniform filtering wastes
capacity by treating both the same. \DFRB{} splits features implicitly
using LIF thresholding: inputs strong enough to cross the firing
threshold produce spikes (the high-frequency component); the unfired
residual retains the low-frequency complement. The temporal split gives
\begin{equation}
    \mathbf{x}_h^{(1)} = \mathrm{LIF}(\mathbf{x}), \qquad
    \mathbf{x}_l^{(1)} = \mathbf{x} - \mathbf{x}_h^{(1)},
    \label{eq:dfrb-split1}
\end{equation}
where $\mathbf{x}\in\mathbb{R}^{T\times C\times H\times W}$ is the
input feature tensor, $\mathbf{x}_h^{(1)}$ is the high-frequency
spike component, and $\mathbf{x}_l^{(1)}$ is the low-frequency
residual. A second branch produces the analogous pair
$\mathbf{x}_h^{(2)},\mathbf{x}_l^{(2)}$ at reduced spatial
resolution. The two high-frequency components are coupled by a
cross-scale gate,
\begin{equation}
    \mathbf{x}_{\mathrm{cross}} =
    \sigma(\gamma)\cdot\mathbf{x}_h^{(1)}\odot\mathbf{x}_h^{(2)},
    \label{eq:dfrb-gate}
\end{equation}
where $\gamma\in\mathbb{R}$ is a learnable scalar initialized to
zero, $\sigma(\cdot)$ is the sigmoid function, and $\odot$ denotes
element-wise multiplication. A learnable per-branch scalar is then
applied before the final summation, so each block starts as an
identity and only opens cross-scale coupling when training finds it
useful.

\noindent\textbf{Spectral--Spatial Hybrid Attention Module (\SHAM{}).}
Following the design in~\cite{ref:fstasnn2025}, this module extends
temporal attention with a 2-D spatial-frequency branch, applied in
series. The temporal branch produces a per-step gate
$\mathbf{g}_t\in\mathbb{R}^T$ by mixing average- and max-pooled spatial
statistics with learnable scalars $\alpha$ and $\beta$ through a
time-axis FC layer, where each element of $\mathbf{g}_t$ rescales the
corresponding time step. The spatial branch computes a 2-D real FFT of
the temporally-averaged feature
$\bar{\mathbf{x}}\!=\!\tfrac{1}{T}\sum_t\mathbf{x}[t]$, reweights the
amplitude $|\mathcal{F}(\bar{\mathbf{x}})|$ via a $1{\times}1$ MLP
while preserving the phase $\angle\mathcal{F}(\bar{\mathbf{x}})$ ---
phase encodes the spatial location of edges and structures, so
modifying it would distort image geometry --- and reconstructs via
inverse FFT. A $7{\times}7$ convolution converts the result into a
spatial attention map $\mathbf{A}_s\in\mathbb{R}^{H\times W}$. The
full module output is
\begin{equation}
    \mathrm{SHAM}(\mathbf{x}) \;=\;
    \mathbf{x} + \sigma(\lambda_s)\cdot
    \mathbf{A}_s \odot \mathbf{g}_t \odot \mathbf{x},
    \label{eq:sham}
\end{equation}
where $\lambda_s\in\mathbb{R}$ is initialized to zero so the
module starts as an identity.

% ---------------------------------------------------------------------
\subsection{Network Scale: AGFM and SSHB}
\label{sec:framework-network}
% ---------------------------------------------------------------------

\noindent\textbf{Adaptive Gated Fusion Module (\AGFM{}).}
Standard U-Nets concatenate encoder and decoder features directly,
which can cause inconsistent fusion when the two paths encode
different cloud--surface mixtures. \AGFM{} replaces the
concatenation with a learnable scalar gate per skip connection,
initialized so that the encoder branch dominates at training start;
the gate gradually opens as training discovers useful decoder
information.

\noindent\textbf{Spectro-Temporal Spike Hyper-Block (\SSHB{}).}
Stage~1 operates at full $H{\times}W$ resolution and dominates the
compute budget, so we trade pixels for time steps,
\begin{equation}
    [T,\,C,\,H,\,W]
    \;\xrightarrow{\mathrm{reshape}}\;
    \Bigl[4T,\,C,\,\tfrac{H}{2},\,\tfrac{W}{2}\Bigr].
    \label{eq:sshb-reshape}
\end{equation}
halving the (Multiplication Addition Computation) MAC count while adding only cheap accumulation
overhead on neuromorphic hardware, where extra time steps cost far less
than a larger spatial map. A  3-D convolution with
temporal stride~4 compresses the sequence back to
$[T,C,\frac{H}{2},\frac{W}{2}]$; a frequency
MLP learned from~\cite{ref:darkir2024} then mixes long-range spectral information
before bilinear upsampling restores $[T,C,H,W]$.

% ---------------------------------------------------------------------
\subsection{Training Objective and Federated Aggregation}
\label{sec:framework-loss}
% ---------------------------------------------------------------------

\noindent\textbf{Loss function.}
The network is trained end-to-end with a Charbonnier--SSIM combined
loss,
\begin{equation}
    \mathcal{L}(\hat{\mathbf{y}},\mathbf{y})
    \;=\;
    \frac{1}{|\Omega|}\sum_{p\in\Omega}
    \sqrt{(\hat{y}_p - y_p)^2 + \varepsilon^2}
    \;+\;
    \lambda\bigl(1 - \mathrm{SSIM}(\hat{\mathbf{y}},\mathbf{y})\bigr),
    \label{eq:loss}
\end{equation}
where $\Omega$ is the set of all pixels in a tile, $\hat{y}_p$ and
$y_p$ are the predicted and ground-truth values at pixel $p$,
$\varepsilon\!=\!10^{-3}$ keeps the first term smooth near zero error,
$\lambda\!=\!0.1$ balances the two terms, and
$\mathrm{SSIM}(\cdot,\cdot)\in[-1,1]$ is the structural similarity
index. The Charbonnier term is robust to large pixel errors under thick
cloud; the SSIM term penalizes structural distortion independently of
pixel magnitude.


%\noindent\textbf{Federated aggregation.}
%The same backbone is trained both centrally and federatedly. In the
%federated regime, each satellite minimizes
%Eq.~\eqref{eq:loss} on its local data, and the constellation pipeline
%of Section~\ref{sec:sysmodel} mediates parameter exchange. Satellites
%on different orbital planes observe different cloud climatologies,
%causing BN statistics to diverge. Averaging BN as in
%FedAvg~\cite{ref:fedavg} mixes these statistics and degrades local
%performance. We follow \textsc{FedBN}~\cite{ref:fedbn}, which keeps BN
%parameters local to each client, and extend it to also keep \TDBN{}
%parameters local --- \TDBN{} folds the firing threshold $V_{\mathrm{th}}$
%into the normalization scale and is equally sensitive to local data
%distribution. Only convolutional weights $\mathbf{W}$, gate scalars
%$\gamma$, and frequency MLP weights are averaged across satellites eac

% =====================================================================
% Section IV: Experiments
% =====================================================================
\section{Experiments}\label{sec:experiments}

% ---------------------------------------------------------------------
\subsection{Experimental Setup}
% ---------------------------------------------------------------------

\noindent\textbf{Datasets.}
CUHK-CR~\cite{ref:cuhkcr} provides paired cloudy/clear Sentinel-2
RGB tiles. CR1 (thin cloud) and CR2 (thick cloud) define the two
source domains used throughout. 

\noindent\textbf{Models.}
Three matched-protocol backbones share the same aggregation pipeline.
\textbf{\OrbitALIF{}-SNN} (ours) uses \FQS-LIF, \TDBN, $T\!=\!4$,
and base width $C\!=\!24$ ($2.30$\,M parameters).
\textbf{\OrbitALIF{}-ANN} replaces every \FQS-LIF with ReLU and
every \TDBN\ with BN, keeping the same dataflow.
\textbf{PlainUNet} is a vanilla U-Net used as the aggregation
baseline. Comparing all three under identical training and
aggregation isolates the effect of the backbone alone.

\noindent\textbf{Training.}
Centralized runs use 400 epochs of AdamW with batch size 4, and gradient clipping at $\|\mathbf{g}\|_2\leq 1$.
Ablation runs use 300 epochs. Federated runs use $R\!=\!200$ rounds
of $E\!=\!2$ local steps under Dirichlet $\alpha\!=\!0.1$ on a
$5{\times}10$ Walker-Star constellation as described in
Section~\ref{sec:sysmodel}.

\noindent\textbf{Energy measurement.}
Per-layer forward hooks record MAC count and input non-zero rate.
ANN energy is charged at $4.6$\,pJ/MAC (Horowitz, $45$\,nm
CMOS~\cite{ref:horowitz2014}); SNN energy is charged at
$77$\,fJ/SOP for neuromorphic substrates such as
Loihi-2~\cite{ref:loihi}, consistent
with~\cite{ref:flsnn,ref:esdnet}.

% ---------------------------------------------------------------------
%\subsection{SNN-Based Onboard Cloud Removal Work}
%\label{sec:exp-feasibility}
% 
%---------------------------------------------------------------------

%As the first time to tackle Onboard Cloud Removal SNN federated learning architecture ，we need to %know if this network can achieve useful restoration quality
%within satellite power constraints. We answer this by placing
%\OrbitALIF{}-SNN alongside its ANN counterpart and the PlainUNet
%baseline under identical centralized training, then measuring both
%quality and energy.
% ---------------------------------------------------------------------
\subsection{SNN-Based Onboard FL Cloud Removal Work}
\label{sec:exp-feasibility}
% ---------------------------------------------------------------------

We deploy all three backbones on the $5{\times}10$ Walker-Star
constellation under $R\!=\!200$ rounds of FedBN$+$Gossip with
$\alpha\!=\!0.1$. Table~\ref{tab:federated} reports the results.
\OrbitALIF{}-ANN reaches $22.679$\,dB, exceeding PlainUNet by
$+0.667$\,dB, confirming the multi-domain architecture brings a
clear advantage even under decentralized training.
\OrbitALIF{}-SNN sits $0.793$\,dB below its ANN counterpart
($21.886$\,dB) — this is the direct cost of spiking sparsity.
The payoff is energy: as Figure~\ref{fig:energy} shows,
\OrbitALIF{}-SNN consumes only $0.287$\,mJ per inference on
neuromorphic hardware, a $72.3\times$ reduction
over the ANN baseline at $20.75$\,mJ. On general-purpose CMOS
the upper bound is $17.14$\,mJ, still a $1.21\times$ reduction.
The MAC-weighted average non-zero rate is $0.826$ on CR1 and
$0.823$ on CR2; per-layer analysis shows $20$--$70\%$ non-zero
activation. At a $100$--$300$\,W satellite power budget, the
$0.793$\,dB spiking cost is a worthwhile trade: \OrbitALIF{}-SNN
remains within $0.126$\,dB of the PlainUNet ANN baseline while
consuming a fraction of its energy, making it the only viable
option for onboard deployment.

\begin{figure}[!t]
  \centering
  \includegraphics[width=\columnwidth]{fig6_energy_bars.png}
  \caption{Per-image estimated inference energy on a $64{\times}64$
  CUHK-CR1 centre patch. ANN baseline at $4.6$\,pJ/MAC.
  \OrbitALIF{} is shown as a dual bound: lower at $77$\,fJ/SOP
  for neuromorphic substrates and upper charging
  every non-zero MAC at the full ANN rate.}
  \label{fig:energy}
\end{figure}

\begin{table}[!t]
\centering
\caption{Decentralized federated cloud-removal on
CUHK-CR1$\cup$CR2 under \textsc{FedBN}$+$Gossip, $R\!=\!200$
rounds, Dirichlet $\alpha\!=\!0.1$. Comm.\ is cumulative
inter-plane traffic over shared (non-BN) parameters.}
\label{tab:federated}
\setlength{\tabcolsep}{4pt}
\footnotesize
\begin{tabular}{lccc}
\toprule
Backbone & PSNR (dB) & SSIM & Comm.\ (MB) \\
\midrule
PlainUNet (baseline)    & 22.012          & 0.6819          & 8{,}420  \\
\OrbitALIF{}-ANN        & \textbf{22.679} & \textbf{0.6966} & 14{,}777 \\
\OrbitALIF{}-SNN (ours) & 21.886          & 0.6725          & 14{,}777 \\
\bottomrule
\end{tabular}
\end{table}

% ---------------------------------------------------------------------
\subsection{Robustness Under Orbital Heterogeneity}
\label{sec:exp-heterogeneity}
% ---------------------------------------------------------------------

Different orbital planes observe different cloud climatologies. To
stress this source-level skew we tighten the Dirichlet concentration
tenfold from $\alpha\!=\!0.1$ to $\alpha\!=\!0.01$, a regime in
which most satellites see only one or two cloud regimes.
\OrbitALIF{}-SNN drops by only $-0.126$\,dB
($21.886\!\to\!21.760$\,dB) under this tightening. The mechanism is
plane-local normalization: under \textsc{FedBN} the BN and \TDBN\
affine parameters and running statistics are excluded from the
inter-plane average (Section~\ref{sec:sysmodel}), so each plane
keeps a normalizer matched to its own climatology and inherits no
statistical bias from planes with different distributions.

Even at the milder $\alpha\!=\!0.1$, the partition difficulty alone
is severe: $30$ of the $50$ satellites hold only the five-tile
floor of the Dirichlet split, so each satellite trains on a narrow
slice of the data. PlainUNet -- the same \textsc{FedBN}+Gossip
protocol but a vanilla U-Net backbone -- loses $2.86$\,dB from
centralized to federated training, illustrating the cost of this
imbalance. \OrbitALIF{} narrows the same gap through the richer
per-satellite representations produced by \DFRB\ and \SHAM, which
are still trained federatedly but extract more from each
satellite's limited tile budget.

% ---------------------------------------------------------------------
\subsection{End-to-End Onboard Federated Cloud Removal}
\label{sec:exp-federated}
% ---------------------------------------------------------------------

Figure~\ref{fig:fedpsnr} plots PSNR against communication rounds
for all three backbones. All three converge smoothly; the per-plane
min--max envelope of \OrbitALIF{}-SNN remains narrow throughout,
confirming that chain-Gossip propagates updates evenly across
orbital planes despite the asymmetric Walker-Star topology.
Both \OrbitALIF{} variants transmit $14{,}777$\,MB — identical
to each other because spiking sparsity reduces on-chip energy
but does not alter the weight tensors exchanged during aggregation.
The SNN therefore gains its $72.3\times$ energy advantage at zero
communication penalty.

\begin{figure}[!t]
  \centering
  \includegraphics[width=\columnwidth]{fig5_federated_curves.png}
  \caption{Federated test PSNR versus communication round under
  \textsc{FedBN}$+$Gossip with $\alpha\!=\!0.1$. Per-plane shading
  shows the min--max envelope across the five orbital planes.}
  \label{fig:fedpsnr}
\end{figure}

% ---------------------------------------------------------------------
\subsection{Centralized Analysis and Ablation}
\label{sec:exp-ablation}
% ---------------------------------------------------------------------

The federated results above establish that \OrbitALIF{} works as an
onboard system. We now explain why, using centralized training to
isolate architectural effects from federation noise.

\noindent\textbf{Centralized upper bound.}
Under full data access, \OrbitALIF{}-ANN reaches $25.374$\,dB on
CR1 and $23.197$\,dB on CR2, exceeding PlainUNet by $+0.502$\,dB
and $+0.157$\,dB. Table~\ref{tab:cr1} places \OrbitALIF{} against
seven published cloud-removal models: it sits within $0.81$\,dB of
DE-MemoryNet ($26.183$\,dB) while using $2.30$\,M parameters and
$15.60$\,G SOPs — the lowest computational cost among all listed
methods. The two stronger entries require $36.80$\,M parameters and
are incompatible with onboard deployment.

\begin{table}[!t]
\centering
\caption{Cloud-removal comparison on CUHK-CR1. Baseline results
cited from~\cite{ref:cuhkcr}. Ops.\ denotes MACs for ANN models
and SOPs for \OrbitALIF{}.}
\label{tab:cr1}
\setlength{\tabcolsep}{4pt}
\footnotesize
\begin{tabular}{lcccc}
\toprule
Method & Params (M) & Ops.\ (G) & PSNR (dB) & SSIM \\
\midrule
SpA-GAN~\cite{ref:cuhkcr}      & 0.22  & 19.40  & 20.999 & 0.5162 \\
AMGAN-CR~\cite{ref:cuhkcr}     & 0.24  & 48.48  & 20.867 & 0.4986 \\
CVAE~\cite{ref:cuhkcr}         & 15.56 & 37.13  & 24.252 & 0.7252 \\
MemoryNet~\cite{ref:cuhkcr}    & 3.64  & 548.65 & 26.073 & 0.7741 \\
MSDA-CR~\cite{ref:cuhkcr}      & 3.91  & 53.45  & 25.435 & 0.7483 \\
DE-MemoryNet~\cite{ref:cuhkcr} & 36.80 & 199.15 & 26.183 & 0.7746 \\
DE-MSDA~\cite{ref:cuhkcr}      & 36.80 & 199.15 & 25.739 & 0.7592 \\
\midrule
\textbf{\OrbitALIF{} (ours)}
  & \textbf{2.30} & \textbf{15.60}
  & \textbf{25.374} & \textbf{0.7728} \\
\bottomrule
\end{tabular}
\end{table}

\noindent\textbf{Component ablation.}
We disable one component at a time and retrain on CR1 for 300
epochs. B1 replaces \SHAM\ with identity; B2 replaces dual-group
\DFRB\ with a single-group residual block; B3 replaces \FQS-LIF\
with binary spikes. Table~\ref{tab:ablation} and
Figure~\ref{fig:ablation} order the drops as \SHAM\
($-0.597$\,dB) $>$ \DFRB\ ($-0.516$\,dB) $>$ \FQS-LIF\
($-0.302$\,dB). B1 and B2 both fall \emph{below} PlainUNet
($24.872$\,dB) while B3 stays above: the system gain is a synergy
across \SHAM\ and \DFRB\ rather than a sum of independent
contributions. \FQS-LIF\ refines rather than load-bears —
sparsity and expressivity are decoupled by design.

\begin{table}[!t]
\centering
\caption{Component ablation on CUHK-CR1 (300 epochs).}
\label{tab:ablation}
\setlength{\tabcolsep}{4pt}
\footnotesize
\begin{tabular}{lcccc}
\toprule
Variant & Params (M) & PSNR (dB) & $\Delta P$ (dB) & SSIM \\
\midrule
Full (A1)            & 2.30 & 25.374 & ---      & 0.7728 \\
$-$\SHAM{} (B1)      & 2.22 & 24.777 & $-$0.597 & 0.7620 \\
$-$\DFRB{} dual (B2) & 1.44 & 24.858 & $-$0.516 & 0.7633 \\
$-$\FQS-LIF{} (B3)   & 2.30 & 25.072 & $-$0.302 & 0.7666 \\
PlainUNet (C2)       & 1.31 & 24.872 & $-$0.502 & 0.7683 \\
\bottomrule
\end{tabular}
\end{table}

\begin{figure}[!t]
  \centering
  \includegraphics[width=\columnwidth]{fig4_ablation_bars-1.png}
  \caption{Component ablation on CUHK-CR1. PSNR (left) and SSIM
  (right) for full \OrbitALIF{} (A1) and ablations B1--B3.}
  \label{fig:ablation}
\end{figure}

% ---------------------------------------------------------------------
\subsection{Qualitative Results}
\label{sec:exp-qualitative}
% ---------------------------------------------------------------------

Figure~\ref{fig:qualitative} shows four CR1 test tiles. \OrbitALIF{}
recovers texture and colour consistency across all scenes, with the
clearest gain along thin-cloud edges where the spatial path of
\DFRB\ is most active. The advantage narrows under fully opaque
cloud, consistent with the smaller CR2 margin.

\begin{figure*}[!t]
  \centering
  \includegraphics[width=\textwidth]{fig3_qualitative_grid-1.png}
  \caption{Qualitative cloud-removal on four CUHK-CR1 test tiles.
  Columns: cloudy input / \OrbitALIF{}-SNN (ours) / PlainUNet /
  ground truth. Per-image PSNR annotated on each cell.}
  \label{fig:qualitative}
\end{figure*}
% =====================================================================
% Section V: Conclusion
% =====================================================================
\section{Conclusion}\label{sec:conclusion}


We presented \OrbitALIF{}, the first framework to perform both
onboard training and onboard inference for optical cloud removal
on a LEO satellite constellation. Rather than compressing an
existing ground-side model, we co-designed the neuron activation,
restoration backbone, and federated aggregation protocol from the
ground up for the four physical constraints of Walker-Star
operation.

The $2.30$\,M-parameter SNN backbone reaches $25.37$\,dB on
CUHK-CR1, sitting within $0.81$\,dB of the strongest published
cloud-removal model while consuming only $0.287$\,mJ per inference
on neuromorphic hardware — a $72.3\times$ reduction over the ANN
baseline. Under decentralized FedBN$+$Gossip with Dirichlet
$\alpha\!=\!0.1$, \OrbitALIF{}-SNN reaches $21.886$\,dB and
degrades by only $0.126$\,dB when source heterogeneity tightens
tenfold, confirming that the adapted FedBN rule absorbs
orbital-plane climatology shifts without re-tuning. The
architecture advantage scales under federation: the multi-domain
backbone outperforms PlainUNet by $+0.667$\,dB in the ANN variant,
larger than the centralized margin, with no additional communication
overhead from the spiking design.

Future work includes closing the $2.86$\,dB centralized-to-federated
gap via momentum aggregation or orbital-geometry-aware mixing
matrices, extending the framework to multi-spectral and SAR-assisted
inputs, and validating the energy bracket with hardware-measured
power on Loihi-2, Speck, or Akida substrates.



\begin{thebibliography}{16}

% --- ordered by first appearance in the body ---

\bibitem{ref:earthobs2}
S.~Voigt \emph{et~al.},
``Global trends in satellite-based emergency mapping,''
\emph{Science}, vol.~353, no.~6296, pp.~247--252, Jul.\ 2016.

\bibitem{ref:earthobs1}
X.~X.~Zhu, D.~Tuia, L.~Mou, G.-S.~Xia, L.~Zhang, F.~Xu, and
F.~Fraundorfer,
``Deep learning in remote sensing: A comprehensive review and list
of resources,''
\emph{IEEE Geosci.\ Remote Sens.\ Mag.}, vol.~5, no.~4, pp.~8--36,
Dec.\ 2017.

\bibitem{ref:meraner2020}
A.~Meraner, P.~Ebel, X.~X.~Zhu, and M.~Schmitt,
``Cloud removal in Sentinel-2 imagery using a deep residual neural
network and SAR-optical data fusion,''
\emph{ISPRS J.\ Photogramm.\ Remote Sens.}, vol.~166, pp.~333--346,
Aug.\ 2020.

\bibitem{ref:ebel2021multisensor}
P.~Ebel, A.~Meraner, M.~Schmitt, and X.~X.~Zhu,
``Multisensor data fusion for cloud removal in global and all-season
Sentinel-2 imagery,''
\emph{IEEE Trans.\ Geosci.\ Remote Sens.}, vol.~59, no.~7,
pp.~5866--5878, Jul.\ 2021.

\bibitem{ref:cuhkcr}
Y.~Sui, M.~Cheng, W.~Wang, Y.~Yang, and L.~Yu,
``Diffusion enhancement for cloud removal in ultra-resolution
remote-sensing imagery,''
\emph{IEEE Trans.\ Geosci.\ Remote Sens.}, vol.~62, 2024.

\bibitem{ref:spacecpn1}
J.~Liu, Y.~Shi, Z.~M.~Fadlullah, and N.~Kato,
``Space-air-ground integrated network: A survey,''
\emph{IEEE Commun.\ Surveys Tuts.}, vol.~20, no.~4, pp.~2714--2741,
4th~Quart.\ 2018.

\bibitem{ref:walker1984}
J.~G.~Walker,
``Satellite constellations,''
\emph{J.\ Brit.\ Interplanetary Soc.}, vol.~37, pp.~559--572, 1984.

\bibitem{ref:fedavg}
H.~B.~McMahan, E.~Moore, D.~Ramage, S.~Hampson, and B.~A.~y~Arcas,
``Communication-efficient learning of deep networks from
decentralized data,''
\emph{Proc.\ Int.\ Conf.\ Artif.\ Intell.\ Statist.\ (AISTATS)},
2017.

\bibitem{ref:fedmega2024}
Y.~Wang \emph{et~al.},
``FedMega: Communication-efficient orbit-aware federated learning
for satellite networks,''
\emph{IEEE J.\ Sel.\ Areas Commun.}, 2024.

\bibitem{ref:dfedsat}
M.~Wu, Y.~Hu, X.~Hou, and X.~Wang,
``DFedSat: Decentralized federated learning for low-Earth-orbit
satellite constellations,''
\emph{arXiv preprint arXiv:2407.05027}, 2024.

\bibitem{ref:loihi}
M.~Davies \emph{et~al.},
``Loihi: A neuromorphic manycore processor with on-chip learning,''
\emph{IEEE Micro}, vol.~38, no.~1, pp.~82--99, Jan./Feb.\ 2018.

\bibitem{ref:esdnet}
H.~Song, M.~Cai, Y.~Cao, and Y.~Xu,
``Exploring the potentials of spiking neural networks for image
deraining,''
\emph{arXiv preprint arXiv:2207.02094}, 2022.

\bibitem{ref:fstasnn2025}
K.~Lai \emph{et~al.},
``FSTA-SNN: Frequency-based spatial--temporal attention for spiking
neural networks,''
\emph{Proc.\ AAAI Conf.\ Artif.\ Intell.}, 2025.

\bibitem{ref:darkir2024}
D.~Feijoo \emph{et~al.},
``DarkIR: Robust low-light image restoration,''
\emph{Proc.\ IEEE/CVF Conf.\ Comput.\ Vis.\ Pattern Recognit.\
(CVPR)}, 2025.

\bibitem{ref:horowitz2014}
M.~Horowitz,
``Computing's energy problem (and what we can do about it),''
\emph{IEEE Int.\ Solid-State Circuits Conf.\ (ISSCC) Dig.\ Tech.\
Papers}, pp.~10--14, Feb.\ 2014.

\bibitem{ref:flsnn}
P.~Yang, T.~Wang \emph{et~al.},
``Brain-inspired decentralized satellite learning in space
computing-power networks,''
\emph{arXiv preprint arXiv:2501.15995}, 2025.

\end{thebibliography}

\end{document}
